\documentclass[11pt]{article}
\usepackage{fontspec,amsmath,amssymb,amsthm,hyperref}
\usepackage[margin=1in]{geometry}
\setlength{\emergencystretch}{3em}
\title{Local smooth drift defect and the future cross term}
\author{LangevinLean mathematical exposition}
\date{}
\begin{document}\maketitle
\noindent Source module: \texttt{ProximalLangevin/SplittingProcess.lean}. Compile this proof document with XeLaTeX or LuaLaTeX. The accompanying HTML pairs the mathematical argument with the exact Lean statements and records the scope of the formalization.


Let a continuous Hilbert-valued reference satisfy \(Y_t-Y_0=-\int_0^t g(Y_s)ds+B_t\), with \(\|g(Y_t)\|\le G\) and \(\|B_t\|\le\sigma\sqrt t\). Only the smooth-part Lipschitz constant enters the local defect.


\section{Integrate the smooth drift increment}


The integral equation implies \(\|Y_t-Y_0\|\le Gt+\sigma\sqrt t\). If \(a\) is \(L\)-Lipschitz and \(0\le t\le h\), then \(\|a(Y_t)-a(Y_0)\|\le L(Gh+\sigma\sqrt h)\). Integrating over an interval of length \(h\) yields

\[\left\|\int_0^h(a(Y_t)-a(Y_0))dt\right\|\le Lh(Gh+\sigma\sqrt h).\]


For the integral solution \(Y_t-Y_0=-\int_0^t g(Y_s)ds+B_t\), assume \(\|g(Y_s)\|\le G\) and \(\|B_t\|\le\nu\sqrt t\) in the Hilbert space under consideration. Then \(\|Y_t-Y_0\|\le tG+\nu\sqrt t\). If only the smooth part \(a\) is \(L\)-Lipschitz, \[\left\|\int_0^h(a(Y_t)-a(Y_0))dt\right\|\le Lh(Gh+\nu\sqrt h).\] The code first bounds the integrand uniformly for \(0\le t\le h\), using monotonicity of \(t\) and \(\sqrt t\), and then applies the interval-integral norm inequality. The Lipschitz constant of the full drift ensures regularity but does not enter this numerical bound.




\textbf{Correspondence with the displayed declarations.}


\begin{itemize}

\item \texttt{splitting\_localDefect\_bound}: A coarse local smooth-drift defect bound uses only the smooth part's own Lipschitz constant and the total reference drift moment.
\end{itemize}


\noindent Lean declarations: \texttt{splitting\_localDefect\_bound}.

\section{Cancel future noise before estimating the remainder}


Subtract the reference equations at \(t,h\) and use adjacent-interval additivity:

\[Y_t-Y_h=\int_t^h g(Y_s)ds-(B_h-B_t).\]

If \(\|b\|\le K\) and \(\langle b,B_h-B_t\rangle=0\), the noise term cancels exactly. Cauchy–Schwarz and the drift bound give \(\langle Y_t-Y_h,b\rangle\le(h-t)GK\). This avoids an order-\(\sqrt{h-t}\) loss.




Subtract the two integral equations and split the integral at \(t\): \[Y_t-Y_h=\int_t^h g(Y_s)ds-(B_h-B_t).\] Pair this identity with \(b\). The explicit orthogonality assumption cancels the future-noise contribution before any norm estimate, leaving \[\langle Y_t-Y_h,b\rangle\le\left\|\int_t^h g(Y_s)ds\right\|\|b\|\le(h-t)GK.\] In the probabilistic application this is an inner product in \(L^2\), so orthogonality means the expected inner product is zero; it is not a samplewise cancellation of Brownian increments.




\textbf{Correspondence with the displayed declarations.}


\begin{itemize}

\item \texttt{reference\_future\_inner\_le}: Future-noise orthogonality turns the reference increment cross term into a drift integral of order \texttt{h-t}, rather than a Brownian norm of order \texttt{sqrt(h-t)}.
\end{itemize}


\noindent Lean declarations: \texttt{reference\_future\_inner\_le}.

\end{document}
